\chapter{PREGLED DOSADAŠNJIH ISTRAŽIVANJA}
\label{ch:pregled}

Zadatak iz ovog rada dodiruje tri područja koja se u literaturi razvijaju odvojeno: upravljanje
u kontaktu, koje se bavi time kako robot smije djelovati na okolinu; mobilnu manipulaciju, koja
se bavi koordinacijom platforme i manipulatora; te interakciju s artikuliranim objektima, koja
se bavi time kako robot doznaje kako se objekt uopće giba. Pregled slijedi tu podjelu, a
završava pristupima temeljenima na učenju, koji zadnjih godina povezuju sva tri.

\section{Klasični pristupi upravljanju u kontaktu}
\label{sec:pregled-klasicni}

Upravljanje robotom u kontaktu s okolinom ne može biti čisto pozicijsko. Čim je vrh alata krut,
a okolina također kruta, mala pogreška položaja proizvodi veliku silu, pa se regulator mora
odnositi i prema odnosu položaja i sile.

Dva su rana odgovora na taj problem oblikovala područje. Hibridno silno-pozicijsko upravljanje
prostor zadatka dijeli na smjerove u kojima se regulira položaj i smjerove u kojima se regulira
sila, prema unaprijed poznatoj matrici izbora izvedenoj iz geometrije zadatka
\cite{raibert1981hybrid}. Impedancijsko upravljanje ne razdvaja te smjerove, nego propisuje
željeni dinamički odnos između pomaka i sile, pa se robot prema okolini ponaša kao opruga
zadane krutosti i prigušenja.

Formalni okvir koji povezuje oba pristupa dao je Khatib upravljanjem u operacijskom prostoru
\cite{khatib1987operational}. Jednadžbe gibanja manipulatora ondje se izražavaju izravno u
prostoru zadatka, pa se sila na vrhu alata preslikava u momente zglobova preko transponiranog
Jakobijana. Isti taj odnos se u ovom radu koristi za procjenu sile iz mjerenja momenata
(poglavlje \ref{ch:procjena}). Kod redundantnih manipulatora taj okvir omogućuje i upravljanje
u nul-prostoru, dakle promjenu konfiguracije zglobova bez pomicanja vrha alata.

Primjenu kartezijskog impedancijskog upravljanja upravo na otvaranje vrata mobilnim
manipulatorom prikazali su Ott i suradnici \cite{ott2005door}. Njihov je rad blizak zadatku iz
ovog rada, s tom razlikom što se oslanja na unaprijed poznat model vrata.

Zajednička slabost svih tih pristupa jest osjetljivost na kvalitetu modela. Ako procijenjena
geometrija ograničenja odstupa od stvarne, regulator tu razliku pretvara u kontaktnu silu, i to
tim izraženije što je krutost veća. Upravo se ta osjetljivost u ovom radu pokazala kao granica
klasičnog pristupa na zakretnim vratima (poglavlje \ref{ch:klasicni}).

\section{Mobilna manipulacija}
\label{sec:pregled-mobilna}

Spajanjem manipulatora s mobilnom platformom uklanja se ograničenje dosega, ali se uvodi
problem koordinacije: platforma i manipulator dijele isti zadatak, a imaju vrlo različitu
dinamiku i točnost. Sustavan pregled pristupa planiranju gibanja za mobilne manipulatore daju
Sandakalum i Ang \cite{sandakalum2022review}.

Postoje dvije osnovne strategije. Hijerarhijski pristup platformu i manipulator rješava
odvojeno, najčešće tako da se platforma najprije dovede na povoljan položaj, a zatim manipulator
izvede zadatak iz mirovanja. Prednost je jednostavnost, a slabost to što se gibanja ne preklapaju
i što loš izbor položaja platforme zadatak može učiniti neizvedivim. Cjelovito upravljanje
platformu i manipulator promatra kao jedan sustav s velikim brojem stupnjeva slobode, čime dobiva
na svestranosti, ali uz znatno složenije upravljanje.

Predstavnik cjelovitog pristupa je perceptivno prediktivno upravljanje temeljeno na modelu, u
kojem se gibanje platforme i manipulatora optimira zajedno, uz model okoline izgrađen iz
senzorskih podataka \cite{pankert2020perceptive}. Takvi pristupi daju glatko i usklađeno gibanje,
ali zahtijevaju model zadatka i računalno su zahtjevni.

Klasični pristup u ovom radu je dijelom cjeloviti, dijelom hijerarhijski,
dok pristup temeljen na učenju koordinaciju rješava jednom naučenom strategijom.

\section{Interakcija s artikuliranim objektima}
\label{sec:pregled-artikulirani}

Vrata, ladice i slični objekti dijele svojstvo da im je gibanje ograničeno na jedan ili nekoliko
stupnjeva slobode, koje robot unaprijed ne poznaje. Prije nego takav objekt pomakne, robot mora
utvrditi kako se on uopće smije gibati.

Sturm i suradnici formulirali su taj problem kao učenje kinematičkih modela artikuliranih
objekata iz opažanja gibanja \cite{sturm2011probabilistic}. Njihov okvir iz zabilježene putanje
odabire model koji ju najbolje objašnjava (rotacijski, prizmatični ili krut) i time
istovremeno prepoznaje tip zgloba i procjenjuje njegove parametre. Isto načelo, prilagodba više
modela istoj putanji, stoji iza postupka procjene ograničenja opisanog u poglavlju
\ref{ch:procjena}.

Noviji radovi taj problem rješavaju tijekom samog izvođenja, bez prethodne faze promatranja.
Rizzi i suradnici predstavili su upravljanje mobilnim manipulatorom temeljeno na uzorkovanju,
otporno na nepoznate parametre artikuliranog objekta \cite{rizzi2023articulated}.

Za ovaj je rad posebno blizak rad Stuedea i suradnika \cite{stuede2019door}, u kojem mobilni
manipulator istog tipa kao u ovom radu otvara vrata i prolazi kroz njih uz kartezijsko
impedancijsko upravljanje, bez poznatog modela vrata, pri čemu se kvaka detektira konvolucijskom
neuronskom mrežom. Taj rad pokazuje i put kojim bi se moglo ukloniti najizraženije
pojednostavljenje ovog rada, a to je oslanjanje na vizualne oznake (potpoglavlje
\ref{sec:perception-limits}).

\section{Pristupi temeljeni na učenju}
\label{sec:pregled-rl}

Podržano učenje na zadatke mobilne manipulacije primjenjuje se jer ne zahtijeva eksplicitan model
zadatka: politika strategiju stječe kroz interakciju s okolinom. Cijena je gubitak uvida u
razloge pojedine odluke i ovisnost o vjernosti simulacije u kojoj se uči.

Za zadatke s kontaktom presudan je odabir prostora akcije. Martín-Martín i suradnici pokazali su
da politika koja izravno zadaje momente ili položaje slabo uči zadatke bogate kontaktom, dok
politika kojoj je akcijski prostor varijabilna impedancija (referentni pomak
zajedno s krutošću regulatora) uči znatno uspješnije \cite{martinmartin2019variable}. Politika
tada ne uči samo kamo se gibati, nego i koliko biti popustljiva u pojedinom trenutku. Taj je
akcijski prostor izravna osnova formulacije u poglavlju \ref{ch:rl}.

Drugi smjer istraživanja bavi se time kako naučiti gibanje platforme koje manipulatoru održava
kinematičku izvedivost. Honerkamp i suradnici najprije su pokazali da se politika platforme može
naučiti tako da prati zadanu putanju vrha alata i drži je unutar dosega
\cite{honerkamp2021feasibility}, a zatim taj pristup proširili na proizvoljna gibanja u
nepoznatoj i promjenjivoj okolini \cite{honerkamp2023n2m2}.

Osim samog upravljanja, učenje se koristi i za oblikovanje sustava. Schneider i suradnici
optimiraju parametre montaže manipulatora na platformu, koristeći naučenu politiku u unutarnjoj
petlji, i pokazuju da izbor montaže sam po sebi bitno mijenja uspješnost otvaranja vrata
\cite{schneider2024codesign}. Nalaz je vrijedan i za ovaj rad, jer je manipulator i ovdje
montiran izvan osi platforme, što zakretom platforme proširuje radni prostor (potpoglavlje
\ref{sec:baza-uci}).

Algoritam korišten u ovom radu je proksimalna optimizacija politike (PPO) \cite{schulman2017ppo},
odabrana zbog stabilnosti pri velikom broju paralelnih okruženja.

\section{Položaj ovog rada}
\label{sec:pregled-pozicija}

Navedeni radovi pojedine dijelove zadatka rješavaju bolje nego ovaj rad, ali ga uglavnom
promatraju odvojeno: ili razvijaju upravljanje uz poznat model objekta, ili uče politiku bez
usporedbe s klasičnom alternativom. Doprinos ovog rada nije novi postupak upravljanja, nego
\emph{izravna usporedba} dvaju pristupa izvedenih za istog robota, istu scenu i ista mjerila
uspješnosti, uključujući dva tipa vrata koja isti zadatak čine različito teškim. Time se
pokušava odgovoriti na pitanje pod kojim uvjetima eksplicitno poznavanje geometrije donosi
prednost, a pod kojima je gubi.